\section{Boundaries and Outlook}\label{sec:boundaries}

The theorem proved here is the scoped two-standard-channel rank theorem for the diagonal signed-reversal layer and its type-D parity sublayer (Section~\ref{sec:typed}).  The separator theorem is separately scoped to $X_n=\BtildeI_n(1,1)$.  The enumeration of positive fixed points and the matching-scheme multiplicities are used as classical background and are not claimed as new.

No full native $B_n$ fingerprint theorem, all-irrep saturation theorem, or general Weyl-group theorem is asserted.  The type-D result is confined to the same two ambient standard channels and the same diagonal cells as the full-layer theorem; no other reflection subgroup or channel is treated.  The finite $n=8$ fingerprint material in Appendix~\ref{app:fullfingerprint} is included because it clarifies the boundary of plausible stronger statements: it records exact finite zero mechanisms and selected finite defect accounting, while leaving the remaining atlas rows open.

The natural next mathematical stage is not to broaden the main theorem by wording, but to decide which finite fingerprint mechanisms admit clean all-$n$ statements.  Until such mechanisms are proved uniformly, they remain appendix evidence rather than theorem statements.
